\section{The Low-Period Spectral Frontier}\label{sec:low-period-frontier}

We prove Theorem~\ref{thm:low-period-frontier} by combining complete orbit enumeration with a compressed
set of exact spectral exclusions. Its domain is the set of infinite operators
\eqref{eq:periodic-operator} whose Hamilton-gauge triangle word \(\tau\) has a least positive period
\(p\leq 16\); the minimized quantity is the squared infinite-volume radius \eqref{eq:squared-periodic-radius},
not a finite-holonomy spectrum. A word displayed in a repeated cell is kept
only for complete certificate accounting and is identified with its primitive
phase by zone folding. The theorem does not extrapolate beyond this domain.

\subsection{Complete phase space}

For each displayed cell length \(1\leq p\leq 16\), consider the legal flux words

\begin{equation}
\label{eq:legal-periodic-flux}
Q \in \{+1,-1\}^p,       \prod_i Q_i=1,
\end{equation}

modulo rotations and reflections. Explicit orbit partition and the independent
Burnside calculation of Appendix~\ref{app:orbit-completeness} give the counts

\begin{equation}
\label{eq:low-period-orbit-counts}
1,2,2,4,4,8,9,18,23,44,63,122,190,362,612,1162,
\end{equation}

whose sum is 2,626. Every periodic phase of primitive \(\tau\) period at most 16
appears in the list when written in its primitive cell. Translation,
reflection, global edge-sign negation, and cell repetition have the spectral
equivalences proved in Lemmas 2.1-2.2.

\subsection{Exact exclusion scheme}

For every representative compute the moments \eqref{eq:moment-and-excess}. Lemma~\ref{lem:moment-barrier} supplies the
valid implication

\begin{equation}
\label{eq:frontier-moment-exclusion}
F_k>0 \Longrightarrow R(Q)>8>\eta.
\end{equation}

An adaptive exact closed-walk calculation through \(F_{64}\) finds a first
positive excess for 2,611 representatives. The first-positive indices range
from 1 to 64; in particular, two near-barrier classes are first detected only
at \(F_{48}\) and \(F_{64}\). This calculation uses integer closed-walk sums, so the
sign of every excess is exact.

Fifteen representatives remain. Eight are displayed even-period repetitions
of the all-negative phase. Equation~\eqref{eq:all-negative-square} proves \(R=8>\eta\) for all of them.
Two are

\begin{equation}
\label{eq:target-zone-folding-rows}
\begin{aligned}
p&=8:  Q=00010001, \\
p&=16: Q=0001000100010001,
\end{aligned}
\end{equation}

where zero denotes negative flux and one positive flux. Lemma~\ref{lem:zone-folding} shows that
the second row is the doubled-cell representation of the first, not a second
phase; Section~\ref{sec:period-eight-edge} gives their common value \(R=\eta\).

The remaining five representatives are

\begin{table}[tbp]
\centering
\caption{Residual low-period competitors requiring endpoint certificates.}
\label{tab:residual-competitors}
\begingroup
\renewcommand{\arraystretch}{1.12}
\begin{tabular}{ll}
\toprule
\(p\) & canonical \(Q\) \\
\midrule
10 & \(0000010001\) \\
12 & \(000000010001\) \\
14 & \(00000000010001\) \\
14 & \(00010010001001\) \\
16 & \(0000000000010001\) \\
\bottomrule
\end{tabular}
\endgroup
\end{table}

For each row choose \(z \in \{+1,-1\}\) and a stored nonzero integer vector \(v\).
Let \(H=H_Q(z)\) and

\begin{equation}
\label{eq:residual-rayleigh-quotient}
r=(v^T H^{2} v)/(v^T v).
\end{equation}

All arithmetic in \eqref{eq:residual-rayleigh-quotient} is integral before the final rational division. The
five certificates first verify \(r>4\). Put

\begin{equation}
u=((r-4)^{2}-10)/2.
\end{equation}

They then verify \(u>0\) and \(u^{2}>5\). Since \(r-4>0\), these inequalities select
the positive square-root branch and give

\begin{equation}
\label{eq:residual-radical-comparison}
\begin{aligned}
(r-4)^{2}&>10+2\sqrt{5}, \\
r&>4+\sqrt{10+2\sqrt{5}}=\eta.
\end{aligned}
\end{equation}

The condition \(r>4\) is indispensable: the two inequalities involving \(u\)
alone can also hold on the lower branch. Rayleigh's principle and \eqref{eq:residual-rayleigh-quotient}-\eqref{eq:residual-radical-comparison}
give \(R(Q)\geq r>\eta\) for each residual representative.

\subsection{Completeness accounting}

The exact partition is

\begin{equation}
\label{eq:frontier-partition}
\begin{aligned}2611&\quad\text{moment-detected representatives},\\
8&\quad\text{repeated all-negative representatives},\\
5&\quad\text{endpoint-certified residual competitors},\\
2&\quad\text{displayed target representations},\\
2626&\quad\text{total representatives}.
\end{aligned}
\end{equation}

The sets in \eqref{eq:frontier-partition} are disjoint, and Appendix~\ref{app:orbit-completeness} proves that the 2,626
representatives are the full legal orbit space. Every non-target phase in the
domain therefore has \(R(Q)>\eta\), while the two target rows are one primitive
period-eight phase with \(R(Q)=\eta\). This proves Theorem~\ref{thm:low-period-frontier}.

The nearest observed competitor occurs at period ten, but numerical ordering
within a non-target period is not part of the theorem. What is proved exactly
is the strict comparison of every competitor with \(\eta\).
